\section{Introduction}
\IEEEPARstart{A}{n} electrocardiograph that has drifted out of specification still produces a
trace. A gain error of ten percent, a low-frequency corner that has moved, or a degraded
common-mode rejection do not stop the instrument; they change the amplitudes and the ST segment
that a clinician or an interpretation algorithm will read. The performance that the instrument
must keep is laid down in IEC~60601-2-25~\cite{iec60601225}, but conformance is verified with
external test equipment, at type testing and at periodic maintenance. Between those moments the
instrument has no way of knowing whether its analog front-end still complies, and the deviations
that appear in use, through aging and environment, are precisely those that the user lacks the
equipment to detect~\cite{chen2025}.

What the instrument can do on its own is limited. Integrated front-ends detect a detached
electrode and provide internal test signals to verify the signal chain at power-up; these signals
are applied inside the device and, by the manufacturer's own account, are not intended for
compliance testing~\cite{ads1298}. The calibration pulse lets an operator see a gross gain error.
Neither covers the discrete
network that surrounds the integrated amplifier (protection and bias resistors, the gain-setting
resistor, the coupling and anti-aliasing filters, the driven-right-leg network, the reference
divider), and neither says whether a deviation matters.

Two research communities have worked on parts of this problem without meeting. The analog test
community checks conformance to specifications, because no generally applicable fault model
exists~\cite{cheng2010}. Since the late 1990s, alternate test has replaced direct specification
measurements by cheap signatures from which pass or fail is inferred, with the explicit aim of
reducing the probability of accepting a bad circuit or rejecting a good
one~\cite{cheng2010,variyam1998spec,variyam1998synth,voorakaranam2007}. This work targets
production test of integrated circuits. The fault-diagnosis community, in contrast, asks which
component has failed and answers with machine-learning classifiers trained on simulated
faults~\cite{dieste2021,dieste2024,dieste2025,yang2020,rajpal2026}. There, a fault is defined by
how far a component deviates from its nominal value, typically by a fixed percentage, and the
circuits are a small set of textbook filters. Whether the deviation takes the circuit out of
specification is not part of the question.

For a medical instrument both questions matter, and a third one appears. The instrument must know
whether it is still fit for use, what has to be repaired if it is not, and whether the problem is
in the circuit at all, because the skin-electrode interface is part of the signal path and varies
by orders of magnitude between subjects and electrode types~\cite{joutsen2024}. A fault definition
based on percentage deviation is of little use here: in the front-ends studied in this paper, more
than \SI{93}{\percent} of the circuits with one component \SI{5}{\percent} off nominal, and more
than half of those with one component \SI{50}{\percent} off, still meet every requirement of the
standard.

The choice of circuit is a further issue. Chen \emph{et al.}~\cite{chen2025} have shown that the
component symmetries of the usual benchmark circuits make the responses to different faults
indistinguishable and limit generalization, even when the faulty components are chosen by
sensitivity analysis, and they redesign the training circuit to remove the symmetries. An
instrument designer cannot do that: the circuit is given. What can be done is to find out in
advance which components are indistinguishable and report the diagnosis at that level.

This paper studies in-service self-test of a single-lead ECG front-end by simulation. The
measurements are restricted to what the instrument can acquire with its own converter and
internal test signals. Two circuits are considered: a realistic one, built around an integrated
instrumentation amplifier with a discrete network (\cA), and a discrete three-op-amp amplifier
that stands for the circuits of the diagnosis literature (\cB). The contributions are as follows.
\begin{enumerate}
\item A two-level self-test scheme that joins specification verification and fault location for a
discrete front-end network. The first level predicts compliance with the requirements of
IEC~60601-2-25 from the self-test measurements; the second locates the cause when the circuit
does not comply. On \cA\ a classifier reaches \SI{0.9}{\percent} of escapes at \SI{7.8}{\percent}
of false rejects, where a limit test on the same measurements gives \SI{22.5}{\percent} and
\SI{20.0}{\percent}.
\item A systematic simulation study, with manufacturing tolerances, of a biopotential front-end
with a realistic architecture: about \num{130000} Monte Carlo cases over two circuits, with hard,
parametric, capacitor-degradation, amplifier and electrode faults, each labeled with the value of
eleven clinical specifications.
\item The separation of electrode problems from circuit faults, using the electrode network of
the standard for gel electrodes and parameters measured on dry
electrodes~\cite{iec60601225,joutsen2024}. Non-compliant circuits are recognized in
\SI{94}{\percent} of the cases and electrode problems in \SI{64}{\percent}; the limit is physical,
since a degraded contact on one subject looks like a normal contact on another, and an open
protection resistor is electrically identical to a detached electrode.
\item A quantified comparison of functional and percentage severity. Training the same detector on
``a component is out of tolerance'' instead of ``the circuit is out of specification'' raises
false rejects from \SI{1.6}{\percent} to \SI{37}{\percent}.
\item An open dataset and code for fault diagnosis of a biomedical front-end, with compliance
labels and the continuous specification values.
\end{enumerate}

Two further results follow from the study. The groups of components that cannot be told apart can
be obtained from the sensitivities of the nominal circuit, before any fault is simulated, and
they explain why component-level diagnosis does not generalize to unseen fault magnitudes while
group-level diagnosis does. And the realistic circuit behaves differently from the benchmark-like
one: it has fewer faults that are both non-compliant and invisible, and it is located more
accurately.

The study is based on simulation only, for two reasons. Labeling a case requires the eleven
tests of the standard on a test bench, and the campaign covers some 600 fault conditions with
200 realizations each, most of them requiring a component to be altered; neither is feasible on
hardware at that scale. And the questions asked here (which faults break compliance, which can
be seen, which can be told apart) have to be answered before a bench experiment can be designed
to test a meaningful subset. Validation on hardware is the next step and is outside this paper.
Section~\ref{sec:related} reviews related work, Section~\ref{sec:methods} describes the circuits,
the fault and measurement models and the evaluation protocol, Section~\ref{sec:results} gives the
results, and Sections~\ref{sec:discussion} and~\ref{sec:conclusion} discuss them and conclude.
